\section{从产品经理训练到 AI 应用创业}

导读强调 Lovart 是垂直 Agent 样本，本章先回到陈冕的产品背景，解释为什么这种选择不是偶然。转写中提到他很早就接触产品方法论，也经历过传统互联网产品经理被“造神”的时代。这个背景很重要，因为 AI 应用创业并不是纯模型工程，它要求创始人理解用户、场景、反馈、增长和产品节奏。Lovart 的设计师场景选择，本质上是产品经理式的场景判断：哪里有高频痛点，哪里有强付费意愿，哪里有模型能力刚刚够用的窗口。

\begin{knowledgebox}{产品经理训练如何迁移到 AI 应用}
\begin{tabularx}{\textwidth}{p{0.22\textwidth}p{0.34\textwidth}X}
\toprule
传统产品能力 & 在 AI 应用中的迁移 & 风险 \\
\midrule
用户洞察 & 找到模型能放大的真实痛点 & 把 demo 当需求。 \\
交互设计 & 把复杂模型能力包装成可用流程 & 过度聊天框化。 \\
增长判断 & 抓住 Magic Moment 的传播窗口 & 增长先于留存。 \\
商业化 & 把体验转成付费理由 & 只做免费爽点。 \\
迭代节奏 & 快速试错和收敛 PMF & 跟随模型变化过度摇摆。 \\
\bottomrule
\end{tabularx}
\end{knowledgebox}

\begin{warningbox}{AI 应用不是“产品经理终于不用懂技术”}
恰恰相反，AI 应用创始人必须同时懂模型能力边界和用户工作流。只懂用户，不知道模型能做什么，会错过窗口；只懂模型，不知道用户为什么付费，会停在 demo。
\end{warningbox}

\begin{knowledgebox}{AI 应用创业的四种能力}
\begin{tabularx}{\textwidth}{p{0.20\textwidth}p{0.34\textwidth}X}
\toprule
能力 & 具体含义 & Lovart 场景里的体现 \\
\midrule
模型判断 & 知道底座模型能否支撑任务 & 设计生成、理解、迭代是否足够稳定。 \\
场景选择 & 找到愿意为 AI 付费的人群 & 设计师和创意团队。 \\
产品包装 & 把复杂能力变成可用界面 & 从 brief 到交付的工作台。 \\
窗口运营 & 抓住模型能力刚够用的时刻 & 从低谷到出圈的时间点。 \\
\bottomrule
\end{tabularx}
\end{knowledgebox}

\subsection{本章小结}

Lovart 的创业逻辑来自产品训练和模型窗口的结合。它不是单纯技术公司，也不是传统 SaaS，而是用 AI 重写一个专业工作流。

